\bookchapter{Capitolo 8: Comunicazione: SMS, chiamate e social}

\chapterintro{La comunicazione è la missione più antica e centrale del telefono. Oggi si è trasformata da semplici chiamate e SMS in un ecosistema complesso che include messaggistica istantanea, social network ed email.}

\booksection{Funzioni avanzate dell'app Telefono}

L'app Telefono dei moderni smartphone è molto più di un semplice tastierino: integra strumenti potenti come registrazione delle chiamate, blocco delle chiamate indesiderate e riconoscimento intelligente.

\booksubsection{Registrazione delle chiamate (attenzione alla conformità legale):}

\begin{itemize}
  \item A cosa serve: registrare il contenuto di chiamate importanti per verbali di lavoro, interviste, promesse dei clienti o come prova in caso di controversie.
  \item Nota legale: in Cina, per usare la registrazione di una chiamata come prova è necessario il consenso dell'altra parte, oppure che serva a tutelare i propri diritti legittimi.
  \item Registrazione unilaterale vs consenso reciproco: sul piano giuridico, una registrazione ottenuta «senza il consenso dell'altra parte» non è automaticamente invalida. Secondo le interpretazioni giudiziarie della Corte Suprema Popolare, purché il metodo di acquisizione non violi gravemente i diritti legittimi altrui, non contravvenga a divieti di legge e non contrasti con l'ordine pubblico e il buon costume, la registrazione può comunque essere presentata come prova. Alcuni tribunali provinciali, però, applicano criteri più severi nelle cause civili per le «registrazioni segrete»; nei processi penali la revisione di legalità è di solito più rigorosa.
  \item Differenze locali: i tribunali di Pechino, Shanghai, Guangdong e altre zone sono relativamente tolleranti verso la registrazione unilaterale come prova civile; le registrazioni che coinvolgono segreti commerciali o privacy personale richiedono invece particolare prudenza ovunque e possono essere escluse per violazione della privacy.
  \item Come si fa: su Android di solito c'è un pulsante «Registra» chiaro nella schermata di chiamata; iOS non offre la registrazione nativa delle chiamate e serve un'app di terze parti o un auricolare Bluetooth con funzione di registrazione.
\end{itemize}

\begin{tcbwarn}{Limiti d'uso delle registrazioni}
\begin{itemize}
  \item Non usare le registrazioni per estorcere denaro, diffonderle o insultare pubblicamente qualcuno.
  \item I contenuti che coinvolgono la privacy di terzi (ad esempio conversazioni altrui registrate per caso) non devono essere diffusi a cuor leggero.
  \item Nelle situazioni importanti conviene dichiarare a voce all'inizio della chiamata «questa chiamata verrà registrata»: è più corretto e la prova ha più valore.
\end{itemize}

\end{tcbwarn}

\booksubsection{Gestione approfondita delle chiamate indesiderate:}

\begin{itemize}
  \item Come funziona: l'app Telefono interroga online un vasto database di numeri e mostra «chiamata indesiderata», «chiamata truffa» o il nome dell'azienda durante la chiamata in arrivo.
  \item Come attivarlo: su Android vai in Telefono > Impostazioni > Blocco chiamate indesiderate; su iOS vai in Impostazioni > Telefono > Blocca e identifica chiamate. Per i percorsi dettagliati di ogni marca e le opzioni consigliate, vedi la tabella comparativa nel Capitolo 13 «Strumenti di sistema utili».
  \item Principali piattaforme di segnalazione: in Cina la marcatura dei numeri è gestita soprattutto da tre fornitori di dati — Tencent Mobile Manager, 360 Mobile Guard e Dianhua Bang (Tapir). Anche gli operatori (China Mobile/Unicom/Telecom) usano spesso questi dati per mostrare il «nome del titolare» e per il blocco delle chiamate indesiderate. I database dei produttori di smartphone (Huawei, Xiaomi, OPPO, vivo) collaborano in gran parte con questi fornitori.
  \item Cosa fare se il tuo numero è stato segnalato per errore: se il tuo numero di lavoro è stato marcato erroneamente come «consegna cibo» «corriere» «chiamata indesiderata», puoi presentare ricorso tramite questi canali: Dianhua Bang: visita la piattaforma di ricorso numeri (haoma.tk.cn) e invia un documento d'identità.
  \item 360: invia il ricorso tramite la funzione «Ricorso numero» nell'app 360 Mobile Guard.
  \item Tencent: invia tramite la pagina «Ricorso numero» della mini-app «Tencent Mobile Manager».
  \item I tre grandi operatori: gli utenti China Mobile possono inviare l'SMS «KTFSR» al 10086 per attivare l'anti-spam; gli utenti Telecom seguano l'account ufficiale «Tianyi Anti-Spam»; gli utenti Unicom impostano tutto in «Unicom Mobile Manager».
\end{itemize}

Anti-spam avanzato: attivando il «blocco chiamate ad alta frequenza» dell'operatore (ad es. il limite chiamate ad alta frequenza di China Mobile) si possono bloccare automaticamente le chiamate vocali AI.

Funzioni di blocco integrate nel telefono: Riconoscimento chiamate / blocco intelligente: marca e blocca automaticamente chiamate indesiderate, truffe e pubblicità in base al database numeri.

Lista nera / lista bianca: metti in blacklist i numeri indesiderati noti; aggiungi a whitelist i contatti e i numeri di corrieri/food delivery per evitare falsi positivi.

Filtro parole chiave SMS: gli SMS che contengono parole come «premio» «prestito» «clicca sul link» finiscono automaticamente nella cartella bloccati.

Silenzia chiamate sconosciute su iOS: le chiamate non in rubrica vengono solo silenziate e inviate alla segreteria, senza disturbarti; utile se ricevi molti spam, ma attento a non perdere chiamate importanti di numeri nuovi.

Integrazione con Non disturbare: di notte o durante le riunioni consenti solo le chiamate dei contatti; con l'attivazione programmata la qualità del sonno migliora subito.

Come integrare i software utili: usa prima il blocco integrato del sistema, poi aggiungi l'app National Anti-Fraud Center e l'anti-spam dell'operatore; per la pulizia puoi scegliere strumenti leggeri come Qingzhuo, ma basta una sola app di sicurezza/pulizia della stessa categoria, vedi Capitolo 13.

\booksubsection{Che cos'è VoLTE/VoNR e come attivarlo:}

\begin{itemize}
  \item VoLTE (Voice over LTE): tecnologia di chiamata vocale sulla rete 4G; durante la chiamata non si torna più a 2G/3G, puoi navigare mentre parli, l'audio è più nitido e la connessione più rapida.
  \item VoNR (Voice over NR): chiamata vocale su 5G, l'evoluzione di VoLTE. Latenza inferiore, qualità audio migliore e supporto per videochiamate HD. Serve uno smartphone 5G + un piano 5G + copertura 5G locale.
  \item Come attivarlo: su iOS in Impostazioni > Rete cellulare > Opzioni dati cellulari > Voce e dati attiva «VoLTE»; su Android cerca l'interruttore «VoLTE chiamate HD» in Impostazioni > Rete mobile > Dati mobile > Avanzate. Di solito l'operatore lo attiva in modo predefinito; se durante una chiamata la rete si disconnette e non navighi, quasi sempre VoLTE non è attivo.
\end{itemize}

\begin{tcbinfo}{Lo sapevi?}
Se durante una chiamata la navigazione rallenta o si interrompe, quasi sempre VoLTE/VoNR non è attivo e il telefono è tornato su 2G/3G. Attivandolo puoi «parlare e guardare video» contemporaneamente.

\end{tcbinfo}

\booksection{SMS e iMessage/RCS}

Gli SMS tradizionali non sono scomparsi: restano un canale affidabile per ricevere codici di verifica, avvisi bancari e comunicazioni ufficiali.

\begin{itemize}
  \item iMessage: messaggi inviati via internet tra dispositivi Apple, con bolle blu. Supporta foto e video di alta qualità, conferme di lettura e altro, completamente gratuito e con crittografia end-to-end.
  \item RCS: protocollo di comunicazione di nuova generazione spinto da operatori globali e Google. Funziona in modo simile a iMessage o WeChat, ma basato sul numero di telefono.
  \item Interoperabilità tra iMessage e RCS: da iOS 18 Apple supporta lo standard RCS: iPhone e Android possono scambiarsi immagini HD, conferme di lettura e indicatori di digitazione, superando il divario delle «bolle verdi». Serve però che entrambi gli operatori abbiano RCS attivo e che il sistema sia abbastanza aggiornato.
\end{itemize}

\booksubsection{Situazione di RCS in Cina:}

\begin{itemize}
  \item China Mobile: ha implementato RCS per prima; in alcune province il servizio «5G Message» supporta messaggi multimediali e schede interattive. Serve un telefono compatibile e l'attivazione in filiale.
  \item China Telecom: offre il servizio tramite «Telecom 5G Message»; alcuni piani lo includono di default, i vecchi utenti devono inviare un SMS di attivazione specifico.
  \item China Unicom: ha distribuito 5G Message nelle principali città, ma la copertura dei dispositivi è irregolare; molti Android devono ancora attivare manualmente «Messaggi avanzati» o «5G Message» nell'app Messaggi.
  \item Esperienza reale: in Cina la diffusione di RCS è più lenta del previsto, soprattutto per la forza dell'ecosistema WeChat e la non uniformità dei dispositivi. La comunicazione quotidiana resta dominata da WeChat; RCS svolge più un ruolo di «canale ufficiale» (avvisi bancari, messaggi della pubblica amministrazione).
\end{itemize}

\begin{tcbwarn}{Sicurezza dei codici di verifica}
\begin{itemize}
  \item Non inoltrare mai i codici di verifica: chi ti chiede di «leggere il codice» o «inoltrare l'SMS» può essere un truffatore, anche se si spaccia per «assistenza clienti» «corriere» o «forze dell'ordine». Il codice di verifica = i tuoi soldi.
  \item Anti-phishing SMS: diffida dei link in messaggi come «carta ETC scaduta» «tessera sanitaria bloccata» «rimborso per pacco smarrito»; le istituzioni legittime non chiedono numero di carta o password via link SMS.
  \item Riconoscimento del numero: gli SMS ufficiali delle banche arrivano di solito da numeri brevi che iniziano con 95 o da 1069 + firma aziendale; qualsiasi «avviso bancario» da un normale cellulare a 11 cifre è quasi sempre una truffa.
  \item Backup a due fattori: non cambiare a cuor leggero il numero legato ai conti importanti; se devi cambiarlo, prima scollega tutti i conti chiave, altrimenti il nuovo titolare potrebbe reimpostare le tue password via SMS.
\end{itemize}

\end{tcbwarn}

\booksection{WeChat/QQ: gestione delle super-app social}

WeChat e QQ sono andati oltre la semplice messaggistica e sono diventati «super-app» che uniscono social, pagamenti e informazioni.

\booksubsection{Backup e migrazione della cronologia chat:}

\begin{enumerate}
  \item Accedi a WeChat per PC o Mac sul computer.
  \item Clicca sul menu in basso a sinistra > [Migrazione e backup] > [Backup e ripristino].
  \item Scegli «Esegui backup della cronologia chat sul computer» e seleziona sul telefono le chat da salvare.
  \item Migrazione su un telefono nuovo: sul vecchio telefono WeChat [Io] > [Impostazioni] > [Generali] > [Migrazione e backup cronologia chat] > [Migra] > [Migra su telefono/tablet], inquadra il codice QR e il gioco è fatto.
\end{enumerate}

\booksubsection{Metodi avanzati per liberare spazio:}

\begin{itemize}
  \item Percorso standard: WeChat [Io] > [Impostazioni] > [Generali] > [Spazio di archiviazione]; puoi cancellare la cache e la cronologia chat in modo mirato per data e tipo di file.
  \item Pulizia dall'interno di una chat: apri una chat > «…» in alto a destra > [Cronologia chat] > [Gestisci spazio]; puoi eliminare immagini, video e file di quella chat mantenendo i testi.
  \item Pulizia di foto/video originali: i file HD originali ricevuti nei gruppi occupano più spazio; in [Gestisci spazio] puoi selezionare «Originali» e cancellarli in blocco.
  \item Disattiva il download automatico: [Impostazioni] > [Generali] > [Foto, video, file e chiamate], spegni «Download automatico» per evitare che su Wi-Fi vengano scaricati silenziosamente molti file di gruppo.
  \item Cache delle mini-app: in [Scopri] > [Mini-app] tieni premuto sulle mini-app poco usate e eliminali per liberare la cache locale.
  \item Backup compresso della cronologia: quando esegui il backup da WeChat per PC, spunta «Comprimi» per ridurre alcuni GB di chat al 30\%-50\% della dimensione originale.
\end{itemize}

\booksubsection{Trucchi di gestione dei gruppi:}

\begin{itemize}
  \item Annuncio del gruppo: il proprietario/gli amministratori lo pubblicano da «…» in alto a destra > [Annuncio del gruppo]; tutti ricevono un avviso prioritario, ideale per comunicazioni importanti. Compare automaticamente ai nuovi membri.
  \item Da fare del gruppo: il proprietario può trasformare un messaggio in evidenza in «Da fare del gruppo»; a tutti i membri appare un indicatore rosso che sparisce al completamento. Utile per iscrizioni a eventi e raccolta di compiti.
  \item Catene di gruppo: digita “\#catena” nella casella di input per avviare una catena, ideale per ordinare cibo o raccogliere adesioni. I risultati della catena vengono riassunti automaticamente in cima alla chat di gruppo, evitando che i messaggi inondino la conversazione.
  \item Live del gruppo: «+» in alto a destra > [Live del gruppo]; puoi avviare una diretta visibile solo ai membri del gruppo, utile per formazione aziendale, riunioni di famiglia o beta test di prodotti.
  \item Comprimi la chat del gruppo: nella lista gruppi scorri il nome verso sinistra > [Comprimi]; i gruppi poco importanti finiscono in «Gruppi compressi» e ti avvisano solo quando ti citano.
  \item Solo proprietario/amministratori possono @tutti: nei gruppi con più di 500 membri i membri normali non possono menzionare tutti, per evitare disturbi.
\end{itemize}

\booksubsection{Gestione fine dei permessi di WeChat Moments:}

\begin{itemize}
  \item Visibilità per gruppi: quando pubblichi, clicca «Chi può vedere» e scegli «Solo alcuni» o «Escludi qualcuno», basandoti sulle etichette preimpostate («colleghi» «famiglia» «clienti») per condividere in modo mirato.
  \item Intervallo di visualizzazione per gli amici: [Io] > [Impostazioni] > [Permessi amici] > [Moments] > [Intervallo di visualizzazione per gli amici], puoi scegliere «ultimi tre giorni/un mese/sei mesi/tutto»; consigliamo «sei mesi» per bilanciare privacy e socialità.
  \item Nascondere un singolo post: tieni premuto su un aggiornamento > [Non vedere i suoi Moments] o [Non seguire più]; l'altra persona non lo saprà.
  \item Gestione della pubblicità in Moments: [Io] > [Impostazioni] > [Info su WeChat] > [Informazioni personali e permessi] > [Gestione pubblicità personalizzata]; puoi disattivare le raccomandazioni basate sul comportamento.
  \item Dieci post per gli estranei: [Permessi amici] > [Moments] > [Consenti agli estranei di vedere dieci Moments]; disattivandolo gli sconosciuti non possono sfogliare i tuoi vecchi post.
\end{itemize}

\booksection{Confronto tra videochiamate}

Ogni strumento di videochiamata ha i suoi punti di forza: scegliere quello giusto per la situazione rende la comunicazione più fluida.

\begin{itemize}
  \item Videochiamate WeChat: copertura capillare, adatto a tutte le età. Ideale per contatti quotidiani con amici e parenti. Fino a 20 persone in video (serve la mini-app «Video di gruppo»); difetto: la qualità cala molto se la rete è scarsa.
  \item FaceTime: esclusivo dell'ecosistema Apple, qualità eccellente e crittografia end-to-end. Da iOS 15 si può invitare utenti Android/Web tramite link, ma l'host deve usare un dispositivo Apple. Ideale per famiglie con dispositivi misti.
  \item Zoom: prima scelta per riunioni internazionali, supporta riunioni gratuite fino a 100 persone per 40 minuti. Funzioni complete: condivisione schermo, sfondo virtuale, sale breakout, adatto a collaborazioni internazionali/da remoto.
  \item Tencent Meeting: prima scelta per l'ufficio in Cina, integrazione profonda con WeChat, si entra scansionando un QR. Fino a 300 partecipanti (versione a pagamento), registrazione su cloud, sottotitoli in tempo reale e verbali automatici; ideale per formazione aziendale e grandi riunioni.
  \item Consigli d'uso: per le conversazioni 1:1 usa WeChat; tra utenti Apple usa FaceTime; per riunioni di lavoro formali usa Zoom o Tencent Meeting (in Cina Tencent, oltre confine Zoom).
\end{itemize}

\booksubsection{Requisiti di banda e consumo di dati:}

\begin{itemize}
  \item Videochiamata WeChat: circa 500 Kbps-1 Mbps di banda; in qualità standard consuma 3-5 MB al minuto, in HD 8-12 MB.
  \item FaceTime: SD circa 1 Mbps, 720p circa 2 Mbps, 1080p oltre 4 Mbps. Circa 5-15 MB al minuto.
  \item Zoom/Tencent Meeting: 720p circa 1,5-3 Mbps, 1080p 3-5 Mbps. Circa 300-800 MB all'ora; l'host di riunioni molte persone consuma di più.
  \item Consiglio: per le videochiamate lunghe usa il Wi-Fi quando possibile; sulla rete mobile abbassa la qualità (disattiva «Video HD» nelle impostazioni dell'app) per risparmiare dati.
\end{itemize}

\booksubsection{Trucchi per migliorare la qualità delle videochiamate:}

\begin{itemize}
  \item Illuminazione: non metterti controfinestra, guarda la sorgente di luce (luce naturale preferibile) per evitare di apparire in controluce.
  \item Inquadratura: tieni la fotocamera leggermente più in alto degli occhi per evitare il doppio mento; in orizzontale reggi il telefono con entrambe le mani per maggiore stabilità.
  \item Audio: usa le cuffie per evitare eco e ronzii; in ambienti silenziosi non attivare l'altoparlante.
  \item Rete: preferisci il Wi-Fi; su 4G/5G avvicinati alla finestra per evitare attenuazione del segnale; nelle riunioni con molte persone puoi spegnere le altre video e lasciare solo chi parla.
  \item Sfondo sfocato: Tencent Meeting e Zoom supportano la sfocatura AI per nascondere ambienti disordinati. Attento a indossare tinte unite, evita colori simili allo sfondo per non sembrare davvero a un green screen.
\end{itemize}

\begin{tcbinfo}{Suggerimento}
Se la rete va a scatti, spegni prima il video e tieni l'audio: è più educato che riagganciare. Prima di riunioni importanti puoi usare «Test rete» di Tencent Meeting per verificare la banda.

\end{tcbinfo}

\booksection{Email: configurazione e uso sul telefono}

\begin{itemize}
  \item Aggiungere un account email: su iOS Impostazioni > Mail > Account > Aggiungi account; su Android di solito c'è l'app «Email».
  \item Suggerimenti: imposta notifiche diverse per ogni account email; usa i gesti scorrevoli per archiviare, eliminare o contrassegnare rapidamente i messaggi.
\end{itemize}

\booksubsection{Differenze tra i protocolli IMAP, POP3 ed Exchange:}

\begin{itemize}
  \item POP3: dopo aver scaricato la posta dal server l'originale viene eliminato per impostazione predefinita e resta solo in locale. Adatto se usi un solo dispositivo e vuoi leggere offline. Difetto: cambiando dispositivo non vedi la cronologia; non è più consigliato.
  \item IMAP: la posta resta sempre sul server e in locale c'è solo la cache. Più dispositivi (telefono/pc/tablet) restano sincronizzati: letto da un lato, letto ovunque. È il protocollo consigliato oggi.
  \item Exchange (incluso ActiveSync): protocollo Microsoft che sincronizza oltre alla posta anche calendario, contatti e attività. Prima scelta per la posta aziendale. Supportato da Outlook, NetEase Enterprise Mail e Tencent Enterprise Mail.
  \item Cosa scegliere: per la posta personale usa IMAP; per quella aziendale preferisci Exchange, poi IMAP; evita POP3 salvo esigenze di archiviazione particolari.
\end{itemize}

\booksubsection{Come configurare la posta aziendale:}

\begin{enumerate}
  \item iOS: [Impostazioni] > [Mail] > [Account] > [Aggiungi account] > scegli «Microsoft Exchange» o «Altro».
  \item Android: nell'app «Email» integrata o in Outlook scegli «Aggiungi account» e inserisci indirizzo email completo e password.
  \item Dati del server (esempio Exchange): l'indirizzo del server è di solito «mail.dominioazienda.com» o «outlook.office365.com»; nel campo dominio inserisci la sigla dell'azienda.
  \item Se compare «Impossibile verificare il server», inserisci manualmente: server in entrata (IMAP) porta 993, crittografia SSL; server in uscita (SMTP) porta 465 o 587, crittografia TLS.
  \item Dopo l'accesso attiva «Push» in [Impostazioni account] per ricevere la posta in tempo reale e disattiva «Anteprima messaggio» per proteggere i segreti commerciali.
\end{enumerate}

\begin{tcbinfo}{Suggerimento}
La posta aziendale di solito richiede che l'amministratore attivi una «password di autorizzazione client» o un «codice di autorizzazione IMAP/SMTP»: non puoi usare la password di accesso web. Se l'accesso fallisce, genera prima il codice di autorizzazione nelle impostazioni webmail e poi inseriscilo sul telefono.

\end{tcbinfo}

\begin{tcbpractice}{Pratica: esegui un backup della cronologia chat di WeChat sul computer}
\begin{enumerate}
  \item Avvia il backup dal computer: apri WeChat per PC > menu in basso a sinistra > [Migrazione e backup] > [Backup e ripristino].
  \item Scegli il backup sul computer: clicca «Esegui backup della cronologia chat sul computer».
  \item Seleziona sul telefono il contenuto del backup e attendi il completamento.
\end{enumerate}

\end{tcbpractice}
